% !TEX program = lualatex
\documentclass[professionalfonts]{beamer}
\newcommand{\slidemode}{1}  % カラー投影モード
\usepackage{beamer_template}
\usepackage{enumerate}

\newcommand{\LectureNum}{11}
\newcommand{\LectureTitle}{複素フーリエ級数}
\newcommand{\CourseName}{応用数学}
\renewcommand{\TermName}{前期}

\begin{document}

\begin{frame}{第11回　複素フーリエ級数}
  \begin{block}{今回の内容}
  複素フーリエ級数
  \end{block}
  \vfill\hfill{\scriptsize \textcolor{gray}{第3章 フーリエ解析　§1}}
\end{frame}

\begin{frame}{定義：複素フーリエ級数（周期 $2\pi$）}
  \begin{block}{}
  \[
    f(x) = \sum_{n=-\infty}^{\infty} c_{n} e^{inx}, \quad
    c_{n} = (1/2\pi) \int_{-\pi}^{\pi} f(x) e^{-inx} dx
  \]
  {\footnotesize \textcolor{gray}{オイラーの公式 $e^{inx} = \cos(nx) + i \sin(nx)$ を利用}}
  \end{block}
\end{frame}

\begin{frame}{一般周期関数の複素フーリエ級数（周期 $2l$）}
  \begin{block}{}
  \[
    f(x) = \sum_{n=-\infty}^{\infty} c_{n} e^{in\pi x/l}, \quad
    c_{n} = (1/2l) \int_{-l}^{l} f(x) e^{-in\pi x/l} dx
  \]
  \end{block}
\end{frame}

\begin{frame}{例題6　（p.\,88）}
  \begin{exampleblock}{問題}
    次の周期 4 の関数 $f(x)$ の複素フーリエ級数を求めよ\\[2pt]
    $f(x) = x + 1  (-2 \leq x < 2),  f(x+4) = f(x)$\\[2pt]
  \end{exampleblock}
\end{frame}

\begin{frame}{例題6　解答}
  $c_{n} = (1/4) \int_{-2}^{2} (x+1) e^{-in\pi x/2} dx$\\[4pt]
  $n \neq 0$ のとき $: e^{in\pi} = (-1)^n$ を用いて計算\\[4pt]
  $c_{n} = 2(-1)^{n+1} / (in\pi)$\\[4pt]
  $n = 0$ のとき $: c_{0} = (1/4)\int_{-2}^{2}(x+1)dx = 1$
  \\[8pt]\textcolor{red}{\textbf{答}}：$f(x) = 1 + \sum_{n\neq0} 2(-1)^{n+1}/(in\pi) \cdot e^{in\pi x/2}$
\end{frame}

\begin{frame}{演習問題}
  \begin{exampleblock}{自分で解いてみよう}
  \begin{description}
    \item[問6] 次の周期 2 の関数 $f(x), g(x)$ の複素フーリエ級数を求めよ
    $(1) f(x) = \begin{cases} 0 & (-1 \leq x < 0) \\ 1 & (0 \leq x < 1) \end{cases},  f(x+2) = f(x)$\\
    $(2) g(x) = \begin{cases} 0 & (-1 \leq x < 0) \\ x & (0 \leq x < 1) \end{cases},  g(x+2) = g(x)$\\
    \item[問7] 次の周期 4 の関数 $f(x), g(x)$ の複素フーリエ級数を求めよ
    $(1) f(x) = \begin{cases} 2 & (-2 \leq x < 0) \\ 0 & (0 \leq x < 2) \end{cases},  f(x+4) = f(x)$\\
    $(2) g(x) = \begin{cases} 0 & (-2 \leq x < 0) \\ 2-x & (0 \leq x < 2) \end{cases},  g(x+4) = g(x)$\\
  \end{description}
  \end{exampleblock}
\end{frame}

\begin{frame}{問6(1)　解答}
  $(1)\; l=1:\; c_{n} = \dfrac{1}{2}\displaystyle\int_{0}^1 e^{-in\pi x}\,dx = \dfrac{1}{2}\!\left[\dfrac{e^{-in\pi x}}{-in\pi}\right]_{0}^{1}$\\[6pt]
  $n \neq 0:\; c_{n} = \dfrac{1-e^{-in\pi}}{2in\pi} = \dfrac{1-(-1)^n}{2in\pi}$\\[6pt]
  $n = 0:\; c_{0} = \dfrac{1}{2}\displaystyle\int_{0}^1 dx = \dfrac{1}{2}$
  \\[8pt]\textcolor{red}{\textbf{答}}：
  $(1)\; f(x) = \dfrac{1}{2} + \displaystyle\sum_{n\neq0} \dfrac{1-(-1)^n}{2in\pi} \cdot e^{in\pi x}$\\
\end{frame}

\begin{frame}{問6(2)　解答}
  $(2)\; l=1:\; c_{n} = \dfrac{1}{2}\displaystyle\int_{0}^1 x\,e^{-in\pi x}\,dx$\\[6pt]
  $n \neq 0$: 部分積分 $\to\; c_{n} = \dfrac{i(-1)^n}{2n\pi} + \dfrac{(-1)^n - 1}{2n^{2}\pi^{2}}$\\[6pt]
  $n = 0:\; c_{0} = \dfrac{1}{2}\displaystyle\int_{0}^1 x\,dx = \dfrac{1}{4}$
  \\[8pt]\textcolor{red}{\textbf{答}}：
  $(2)\; g(x) = \dfrac{1}{4} + \displaystyle\sum_{n\neq0} \!\left[\dfrac{i(-1)^n}{2n\pi} + \dfrac{(-1)^n-1}{2n^{2}\pi^{2}}\right] e^{in\pi x}$\\
\end{frame}

\begin{frame}{問7(1)　解答}
  $(1)\; l=2:\; c_{n} = \dfrac{1}{4}\displaystyle\int_{-2}^{0} 2\,e^{-in\pi x/2}\,dx$\\[6pt]
  $n \neq 0:\; c_{n} = \dfrac{1}{2}\!\left[\dfrac{e^{-in\pi x/2}}{-in\pi/2}\right]_{-2}^{0} = \dfrac{e^{in\pi}-1}{in\pi} = \dfrac{(-1)^n-1}{in\pi}$\\[6pt]
  $n = 0:\; c_{0} = \dfrac{1}{4}\displaystyle\int_{-2}^{0}2\,dx = 1$
  \\[8pt]\textcolor{red}{\textbf{答}}：
  $(1)\; f(x) = 1 + \displaystyle\sum_{n\neq0} \dfrac{(-1)^n-1}{in\pi} \cdot e^{in\pi x/2}$\\
\end{frame}

\begin{frame}{問7(2)　解答}
  $(2)\; l=2:\; c_{n} = \dfrac{1}{4}\displaystyle\int_{0}^2 (2-x)\,e^{-in\pi x/2}\,dx$\\[6pt]
  $n \neq 0$: 部分積分 $\to\; c_{n} = \dfrac{1}{in\pi} + \dfrac{1-(-1)^n}{n^{2}\pi^{2}}$\\[6pt]
  $n = 0:\; c_{0} = \dfrac{1}{4}\displaystyle\int_{0}^2(2-x)\,dx = \dfrac{1}{4}\!\left[2x-\dfrac{x^{2}}{2}\right]_{0}^2 = \dfrac{1}{2}$
  \\[8pt]\textcolor{red}{\textbf{答}}：
  $(2)\; g(x) = \dfrac{1}{2} + \displaystyle\sum_{n\neq0} \!\left[\dfrac{1}{in\pi} + \dfrac{1-(-1)^n}{n^{2}\pi^{2}}\right] e^{in\pi x/2}$\\
\end{frame}

\end{document}
